% ============================================================
%  Academic Résumé — inspired by the rainbow-beamer theme
%  Requires XeLaTeX or LuaLaTeX
%  Standard class: article
% ============================================================
\documentclass[11pt, a4paper]{article}

% ── Engine check ────────────────────────────────────────────
\usepackage{iftex}
\RequireLuaTeX % or swap to \RequireLuaTeX

% ── Page geometry ───────────────────────────────────────────
\usepackage[
  top    = 2.0cm,
  bottom = 2.0cm,
  left   = 1.4cm,
  right  = 1.4cm,
]{geometry}

% ── Fonts (mirrors beamerfontthemerainbow.sty) ───────────────
% Main: Yanone Kaffeesatz (display / headings feel)
% Sans: Andika            (body text)
% Mono: Iosevka           (code / technical details)
%
% If these fonts are not installed system-wide, swap the names
% for any similar alternatives, e.g.
%   main → "Cabin" or "Raleway"
%   sans → "Source Sans Pro"
%   mono → "Fira Mono"
\usepackage{fontspec}
\usepackage[warnings-off={mathtools-colon,mathtools-overbracket}]{unicode-math}

\setmainfont[
  Path = {fonts/Yanone_Kaffeesatz/},
  UprightFeatures={
    LetterSpace=4.2,
    Weight=400,
  },
  BoldFont            =   *,
  BoldFeatures={
    LetterSpace=4.2,
    Weight=900,
  },
  ItalicFont = *,
  ItalicFeatures={
      LetterSpace=4.2,
      Slant=-5
    }
]{YanoneKaffeesatz-VariableFont_wght.ttf}

\setsansfont[
  Path           = {fonts/Andika/},
  Extension      = .ttf,
  Ligatures      = TeX,
  UprightFont    = *-Regular,
  BoldFont       = *-Bold,
  ItalicFont     = *-Italic,
  BoldItalicFont = *-BoldItalic
]{Andika}

% \setmonofont[
%   UprightFont    = *-Regular,
%   ItalicFont     = *-Italic,
%   BoldFont       = *-Bold,
%   BoldItalicFont = *-BoldItalic,
% ]{Iosevka}

% Use the sans font for body text (matches the beamer theme's
% Andika-as-body approach; swap \rmfamily if you prefer serif).
\renewcommand{\familydefault}{\sfdefault}

% ── Colors (from beamercolorthemerainbow.sty) ───────────────
\usepackage[dvipsnames]{xcolor}

\definecolor{cvgreen}     {HTML}{89B374}
\definecolor{cvlightgreen}{HTML}{B5C266}
\definecolor{cvwhite}     {HTML}{E4E4E4}
\definecolor{cvdarkblue}  {HTML}{387490}
\definecolor{cvnightblue} {HTML}{303042}
\definecolor{cvblue}      {HTML}{01ACB6}
\definecolor{cvblack}     {HTML}{2B2B2D}
\definecolor{cvviolet}    {HTML}{837694}
\definecolor{cvpurple}    {HTML}{6e2c5a}
\definecolor{cvred}       {HTML}{F2503B}
\definecolor{cvyellow}    {HTML}{FEBE28}
\definecolor{cvdarkgray}  {HTML}{655a62}
\definecolor{cvpink}      {HTML}{f67598}
\definecolor{cvorange}    {HTML}{e9a02f}
\definecolor{cvmint}      {HTML}{3eb9a5}

% ── Packages ────────────────────────────────────────────────
\usepackage{tikz}
\usetikzlibrary{calc}
\usepackage{hyperref}
\usepackage{titlesec}
\usepackage{mdframed}
\usepackage{enumitem}
\usepackage{tabularx}
\usepackage{array}
\usepackage{microtype}
\usepackage{fontawesome5}
\usepackage[skip=4pt, indent=0pt]{parskip}
\usepackage{etoolbox}

\hypersetup{
  colorlinks = true,
  urlcolor   = cvdarkblue,
  linkcolor  = cvdarkblue,
}

% ── No page numbers by default ──────────────────────────────
\pagestyle{empty}

% Main text color (after all packages are loaded)
\AtBeginDocument{\color{cvblack}}

% ── Rainbow rule (mirrors the gradient used in the theme) ───
% Declared once in the preamble after tikz is loaded.
% Uses explicit RGB so it works without pgf color name lookup at
% shading-declaration time (color names aren't resolved yet then).
\pgfdeclareverticalshading{rainbowcv}{100bp}{%
  rgb(0bp)=(0.949,0.314,0.231);%   cvred    #F2503B
  rgb(14bp)=(0.996,0.745,0.157);%  cvyellow #FEBE28
  rgb(28bp)=(0.537,0.702,0.455);%  cvgreen  #89B374
  rgb(42bp)=(0.004,0.675,0.714);%  cvblue   #01ACB6
  rgb(56bp)=(0.220,0.455,0.565);%  cvdarkblue #387490
  rgb(70bp)=(0.514,0.463,0.580);%  cvviolet #837694
  rgb(84bp)=(0.431,0.173,0.353);%  cvpurple #6e2c5a
  rgb(100bp)=(0.431,0.173,0.353)%  cvpurple
}

% Usage: \rainbowrule  or  \rainbowrule[3pt]
% Drawn inline with zero surrounding space; titlesec after-code
% adds its own separation via \titlespacing.
\newcommand{\rainbowrule}[1][2pt]{%
  \noindent\textcolor{cvblue}{\rule{\linewidth}{#1}}%
}

% ── Section heading style ───────────────────────────────────
\newfontfamily\headingfont[
  Path = {fonts/Yanone_Kaffeesatz/},
  UprightFeatures = {LetterSpace=5, Weight=700},
]{YanoneKaffeesatz-VariableFont_wght.ttf}

% Section: title on the left, rainbow rule fills the rest of the line.
% \titleformat's "before-code" (5th arg) runs before the title text,
% and "after-code" (6th arg in []) runs after. We put the rule in the
% after-code as a negative-vspace trick to pull it up to the title baseline,
% but that gives a separate line. Instead we abuse the label (3rd arg) to
% be empty and use the format (2nd arg) + a \dotfill-style construction
% by redefining \section via a wrapper macro.
%
% Cleanest approach: patch \section to a custom macro that typesets
% title + leaders in a single \hbox.
\titleformat{\section}{\LARGE\bfseries\color{cvnightblue}\headingfont}{}{}{}[]
\titlespacing*{\section}{0pt}{16pt}{4pt}

\let\oldsection\section
\renewcommand{\section}[1]{%
  \oldsection*{%
    \begin{tabularx}{\linewidth}{@{} l @{\hspace{6pt}} X @{}}
      #1 &
      \raisebox{0.4ex}{\rainbowrule[1.5pt]}
    \end{tabularx}%
  }%
  \addcontentsline{toc}{section}{#1}%
}

\titlespacing*{\section}{0pt}{16pt}{4pt}

\titleformat{\subsection}
  {\large\bfseries\color{cvdarkblue}\sffamily}
  {}
  {0pt}
  {}

\titlespacing*{\subsection}{0pt}{6pt}{4pt}

% ── CV entry macros ─────────────────────────────────────────

% Uniform inter-entry gap used by all CV macros.
\newlength{\cvskip}
\setlength{\cvskip}{10pt}

% ── Timeline environment ─────────────────────────────────────
% Strategy: savebox pre-measurement.
%
% For each entry:
%   1. Typeset the right-column content into \cvtlbox (a savebox).
%   2. Measure its total height = \ht\cvtlbox + \dp\cvtlbox.
%   3. Typeset a two-column tabular:
%        left  col: TikZ picture whose bar height = measured content height
%        right col: \usebox{\cvtlbox}
%
% The savebox is typeset at the correct text column width, so no overflow.
% Works across page breaks because every entry is self-contained.
%
% Bar extends from the dot downward to the bottom of the content box,
% plus \cvskip, so consecutive bars connect perfectly.
% The last entry in a cvtimeline gets a shorter bar (just below the dot).

\newlength{\cvtldot}    \setlength{\cvtldot}{6pt}     % dot diameter
\newlength{\cvtlindent} \setlength{\cvtlindent}{8pt}   % bar x-offset from left margin
\newlength{\cvtlgap}    \setlength{\cvtlgap}{10pt}     % gap between bar col and content
\newlength{\cvtlcol}    % width of bar column (bar + gap)
\setlength{\cvtlcol}{\dimexpr\cvtlindent+\cvtlgap\relax}
\newlength{\cvtlcontentw} % computed content column width (set at start of each entry)

\newlength{\cvtlentryht} % height of current entry content box
\newlength{\cvtldotyoff} % y-offset from top of box to dot centre (positive = downward)
\newlength{\cvtlbary}    % y-coordinate of bar bottom (from top of box, positive = downward)
\newsavebox{\cvtlbox}    % savebox for entry content

% Track whether this is the last entry in a timeline.
% The environment end-code patches the last entry's bar.
% We use a simple boolean toggled by the environment.
\newif\ifcvtllast
\cvtllastfalse

\newenvironment{cvtimeline}{%
  \par\vspace{2pt}%
  \setlength{\parskip}{0pt}%  suppress parskip so only \cvskip controls gaps
  \cvtllastfalse%
}{%
  \par\vspace{2pt}%
}

% Internal macro that typesets a single timeline entry.
% #1=date  #2=title  #3=org  #4=unused  #5=description
% #6=boolean: 1 = last entry in group (bar to bottom of content only, no extension)
%             0 = non-last (bar extends to meet the next entry's dot)
\newcommand{\cvtlentry}[6]{%
  \setlength{\fboxsep}{2.5pt}%
  \setlength{\cvtlcontentw}{\dimexpr\linewidth-\cvtlindent-\cvtlgap\relax}%
  % ── 1. Typeset content into savebox at the correct width ────
  \savebox{\cvtlbox}{%
    \begin{minipage}[t]{\cvtlcontentw}%
      {\colorbox{cvblue!12}{%
        \fontsize{8}{9}\selectfont\sffamily\color{cvblue!80!black} #1}}%
      \smallskip\newline
      {\normalsize\bfseries\color{cvnightblue} #2}%
      \ifblank{#3}{}{\newline{\small\itshape\color{cvdarkblue} #3}}%
      \ifblank{#5}{}{%
        \newline\vspace{-4pt}%
        \begin{mdframed}[
          topline=false,bottomline=false,rightline=false,
          linewidth=2pt,linecolor=cvdarkblue!30,
          innerleftmargin=6pt,innerrightmargin=0pt,
          innertopmargin=1pt,innerbottommargin=1pt,
          skipabove=0pt,skipbelow=0pt,backgroundcolor=white]%
          {\footnotesize\color{cvdarkgray} #5}%
        \end{mdframed}%
      }%
    \end{minipage}%
  }%
  % ── 2. Geometry ─────────────────────────────────────────────
  % \ht\cvtlbox = height above the minipage's first-line baseline
  %             = top-of-box to first-line-baseline
  % \dp\cvtlbox = depth below that baseline
  % Total box height = \ht + \dp
  %
  % The dot should sit vertically centred in the date badge colorbox.
  % The colorbox is the first thing in the minipage; its centre is at:
  %   y_dot = \ht\cvtlbox - \fboxsep - (badge_height/2)
  % badge_height = font_size*leading + 2*\fboxsep = 9pt + 5pt = 14pt approx.
  % Simpler and robust: centre of colorbox = \ht\cvtlbox - \fboxsep - 4.5pt
  % We store y_dot as a length measured from the top of the box (positive downward).
  %
  % In TikZ we work with (0,0) at the TOP of the box (achieved via \raisebox),
  % so dot y-coordinate = -\cvtldotyoff (negative = downward from top).
  \setlength{\cvtldotyoff}{\dimexpr\fboxsep + 4pt\relax}%  ≈ centre of 8/9pt badge
  %
  % Bar length:
  %   non-last: from dot centre to bottom of next entry's dot
  %             = (box_height - dot_yoff) + \cvskip + dot_yoff_next
  %             = \dp\cvtlbox + \ht\cvtlbox - \cvtldotyoff + \cvskip + \cvtldotyoff
  %             = total_box_height + \cvskip
  %   last:     from dot centre to bottom of content box
  %             = total_box_height - \cvtldotyoff
  \ifnum#6=0\relax
    \setlength{\cvtlentryht}{\dimexpr\ht\cvtlbox+\dp\cvtlbox+\cvskip\relax}%
  \else
    \setlength{\cvtlentryht}{\dimexpr\ht\cvtlbox+\dp\cvtlbox-\cvtldotyoff\relax}%
  \fi
  % ── 3. Typeset ──────────────────────────────────────────────
  % Pre-compute the bar's bottom y-coordinate (from top of box, positive downward).
  \setlength{\cvtlbary}{\dimexpr\cvtldotyoff+\cvtlentryht\relax}%
  %
  % Layout: zero-width gutter on the left containing the bar+dot overlay,
  % then the content minipage.  No tabular — avoids p{} leading issues.
  %
  % \raisebox{\ht\cvtlbox}[0pt][0pt] shifts (0,0) of the tikzpicture to the
  % TOP of the content box, so all y-offsets are negative (downward).
  \noindent
  \hspace{\cvtlindent}%
  \raisebox{\ht\cvtlbox}[0pt][0pt]{%
    \makebox[0pt][c]{%
      \begin{tikzpicture}[baseline=0pt, inner sep=0pt, outer sep=0pt]
        % bar first (dot will be drawn on top)
        \draw[cvdarkgray!40, line width=1pt]
          (0pt, -\cvtldotyoff) -- (0pt, -\cvtlbary);
        % dot
        \node[circle, fill=cvblue, minimum size=\cvtldot,
              inner sep=0pt, outer sep=0pt]
          at (0pt, -\cvtldotyoff) {};
      \end{tikzpicture}%
    }%
  }%
  \hspace{\cvtlgap}\usebox{\cvtlbox}%
  \par\vspace{\cvskip}%
}

% Public \cventry: forwards to \cvtlentry with last=0.
% The last entry must be wrapped with \cvlastentry instead,
% or users can use \cventry and accept a trailing stub bar.
\newcommand{\cventry}[5]{\cvtlentry{#1}{#2}{#3}{#4}{#5}{0}}
\newcommand{\cvlastentry}[5]{\cvtlentry{#1}{#2}{#3}{#4}{#5}{1}}

% \cvpub{authors}{title}{venue/journal}{year}{url (optional)}
\newcommand{\cvpub}[5]{%
  \noindent
  {\small
    \textcolor{cvdarkgray}{#1}. \textbf{\textcolor{cvnightblue}{#2}}.
    \textit{\textcolor{cvdarkblue}{#3}}, #4.%
    \ifblank{#5}{}{\enspace\href{#5}{\textcolor{cvblue}{\faLink}}}
  }\vspace{\cvskip}\par\noindent%
}

% \cvskillblock{icon}{Category}{tag, tag, ...}
% Renders a category label with a FontAwesome icon and \cvtag pills.
% Use inside the cvskills environment.
% Note: uses a minipage so tabularx doesn't mis-scan the cell content.
\newcommand{\cvskillblock}[3]{%
  \begin{minipage}[t]{\linewidth}%
    {\fontsize{8}{9}\selectfont\sffamily\color{cvdarkgray}%
     \textcolor{cvdarkblue}{#1}\enspace\MakeUppercase{#2}}%
    \vspace{3pt}\par
    \setlength{\fboxsep}{3pt}\raggedright\sloppy
    #3\par
  \end{minipage}%
}

% cvskills environment: tinted box, two-column grid via two minipages side by side.
\newlength{\cvskillcolw}
\newenvironment{cvskills}{%
  \begin{mdframed}[
    topline=false, bottomline=false, leftline=false, rightline=false,
    backgroundcolor=cvdarkblue!4,
    innerleftmargin=10pt, innerrightmargin=10pt,
    innertopmargin=8pt, innerbottommargin=8pt,
    skipabove=0pt, skipbelow=0pt]%
  \setlength{\cvskillcolw}{\dimexpr0.5\linewidth-6pt\relax}%
}{%
  \end{mdframed}%
}

% Contact icon helpers
\newcommand{\cvcontact}[2]{\textcolor{cvdarkblue}{#1}\enspace\href{#2}{#2}}
\newcommand{\cvemail}[1]{\textcolor{cvdarkblue}{\faEnvelope}\enspace\href{mailto:#1}{#1}}
\newcommand{\cvphone}[1]{\textcolor{cvdarkblue}{\faPhone*}\enspace #1}
\newcommand{\cvweb}[2]{\textcolor{cvdarkblue}{\faGlobe}\enspace\href{#1}{#2}}
\newcommand{\cvgithub}[1]{\textcolor{cvdarkblue}{\faGithub}\enspace\href{https://github.com/#1}{github.com/#1}}
\newcommand{\cvlinkedin}[1]{\textcolor{cvdarkblue}{\faLinkedin}\enspace\href{https://linkedin.com/in/#1}{in/#1}}
\newcommand{\cvorcid}[1]{\textcolor{cvdarkblue}{\faOrcid}\enspace\href{https://orcid.org/#1}{#1}}

% ── Pull-quote block (Research Interests) ───────────────────
% Left accent bar in cvdarkblue, very light tint background.
\newmdenv[
  topline    = false,
  bottomline = false,
  rightline  = false,
  leftline   = false,
  backgroundcolor = cvdarkblue!4,
  innerleftmargin  = 10pt,
  innerrightmargin = 8pt,
  innertopmargin   = 6pt,
  innerbottommargin= 6pt,
  skipabove  = 0pt,
  skipbelow  = 0pt,
]{pullquote}

% \cvtag{text} — small pill for interest keywords, wraps naturally.
% Must be used inside \begin{cvtags}...\end{cvtags}.
\newenvironment{cvtags}{%
  \begin{sloppypar}\setlength{\fboxsep}{3pt}\raggedright%
}{%
  \end{sloppypar}%
}
\newcommand{\cvtag}[1]{%
  {\fontsize{8}{10}\selectfont\sffamily
   \colorbox{cvblue!12}{\textcolor{cvblue!80!black}{#1}}}%
  \hspace{4pt plus 2pt minus 1pt}\linebreak[0]%
}

% ── List style ───────────────────────────────────────────────
\setlist[itemize]{
  label       = \textcolor{cvdarkblue}{\bullet},
  nosep,
  leftmargin  = 1.4em,
  topsep      = 2pt,
  itemsep     = 1pt,
}

% ============================================================
%  DOCUMENT
% ============================================================
\begin{document}

% ── Header ──────────────────────────────────────────────────
\begin{center}
	{\Huge\bfseries\headingfont%
		\textcolor{cvnightblue}{Ghilain} \textcolor{cvblue}{Bergeron}}

	\vspace{4pt}
	{\normalsize\color{cvdarkgray} Ph.D.\ Student in Computer Science --- Formal Methods}

	\vspace{1pt}
	\begin{minipage}{0.9\linewidth}\centering\rainbowrule[1pt]\end{minipage}

	{\small
	\cvemail{ghilain.bergeron@inria.fr}
	\quad\textcolor{cvdarkgray}{|}\quad
	\cvgithub{mesabloo}
	\quad\textcolor{cvdarkgray}{|}\quad
	\cvorcid{0009-0008-2814-3480}
	% \quad\textcolor{cvdarkgray}{|}\quad
	\cvweb{https://g-bergeron.github.io}{https://g-bergeron.github.io}
	\quad\textcolor{cvdarkgray}{|}\quad
	\cvlinkedin{g-bergeron}
	}
\end{center}

\vspace{8pt}

% ── Research interests ──────────────────────────────────────
\section{\faFlask\enspace Research Interests}

\begin{pullquote}
	% I am a Ph.D.\ student at Inria Nancy in the VeriDis team, advised by
	% Prof.\ Horatiu Cirstea and Prof.\ Stephan Merz.
	My research focuses on \textit{compiler correctness}, \textit{programming language semantics},
	and \textit{formal software verification}, with an emphasis on mechanised proofs in the
	Lean~4 proof assistant.
	I am currently working on a formally verified compiler for Distributed PlusCal targeting Go.

	\vspace{6pt}
	\begin{cvtags}
		\cvtag{Formal methods}
		\cvtag{Compiler correctness}
		\cvtag{Lean 4}
		\cvtag{PL semantics}
		\cvtag{Proof assistants}
		\cvtag{Type theory}
	\end{cvtags}
\end{pullquote}


% ── Skills ──────────────────────────────────────────────────
\section{\faStar\enspace Skills}

\begin{cvskills}
	\noindent
	\begin{minipage}[t]{\cvskillcolw}
		\cvskillblock{\faCertificate}{Proof assistants}{%
			\cvtag{Lean 4}\cvtag{Isabelle/HOL}\cvtag{Coq/Rocq}\cvtag{Agda}}
	\end{minipage}%
	\hspace{12pt}%
	\begin{minipage}[t]{\cvskillcolw}
		\cvskillblock{\faCode}{Programming}{%
			\cvtag{Haskell}\cvtag{Go}\cvtag{Python}\cvtag{Rust}\cvtag{C/C++}}
	\end{minipage}

	\vspace{8pt}

	\noindent
	\begin{minipage}[t]{\cvskillcolw}
		\cvskillblock{\faWrench}{Tools}{%
			\cvtag{\LaTeX}\cvtag{Git}\cvtag{Nix}}
	\end{minipage}%
	\hspace{12pt}%
	\begin{minipage}[t]{\cvskillcolw}
		\cvskillblock{\faComments}{Languages}{%
		\cvtag{French -- {\footnotesize\textit{native}}}\cvtag{English -- {\footnotesize\textit{fluent}}}}
	\end{minipage}
\end{cvskills}

% ── Education ───────────────────────────────────────────────
\section{\faGraduationCap\enspace Education}

\begin{cvtimeline}
	\cventry{2023 -- present}
	{Ph.D.\ in Computer Science}
	{University of Lorraine --- Inria Nancy, VeriDis Team}
	{}
	{\begin{itemize}
			\item Thesis topic: Generation of distributed programs from formal specifications.
			\item Advisors: \href{https://members.loria.fr/HCirstea/}{Horatiu Cirstea} and \href{https://members.loria.fr/Stephan.Merz/}{Stephan Merz}.
		\end{itemize}}

	\cventry{2021 -- 2023}
	{M.Sc.\ in Computer Science}
	{University of Lorraine --- Faculté des Sciences et Technologies}
	{}
	{\begin{itemize}
			\item Specialty: Software Engineering, track FM2S (Formal Methods for Safe Software).
			\item Ranked 1\textsuperscript{st} out of 21 students.
		\end{itemize}}

	\cvlastentry{2018 -- 2021}
	{B.Sc.\ in Computer Science}
	{University of Lorraine --- IUT Nancy-Charlemagne \& Faculté des Sciences et Technologies}
	{}
	{}
\end{cvtimeline}

% ── Research experience ─────────────────────────────────────
\section{\faMicroscope\enspace Research Experience}

\begin{cvtimeline}
	\cventry{Mar -- Aug 2023}
	{Research Internship}
	{University of Lorraine --- Inria Nancy, VeriDis Team}
	{}
	{\begin{itemize}
			\item Topic: Formalization of a splitting framework for saturation-based calculi in Isabelle/HOL.
			\item Advisor: \href{https://members.loria.fr/sophie.tourret/}{Sophie Tourret}
		\end{itemize}}

	\cventry{Jun -- Aug 2021}
	{Research Internship}
	{University of Lorraine --- Inria Nancy, MFX Team}
	{}
	{\begin{itemize}
			\item Topic: Development of a formal verification pipeline in C++ for
			      \href{https://github.com/sylefeb/Silice}{Silice},
			      a hardware description language, including verification primitives.
			\item Advisor: \href{https://www.antexel.com/sylefeb/research}{Sylvain Lefebvre}
		\end{itemize}}

	\cvlastentry{Apr -- Jun 2020}
	{Research Internship}
	{University of Lorraine --- Inria Nancy, MFX Team}
	{}
	{\begin{itemize}
			\item Development of a tool in C++ for automatic parameterization of 2D-modelling scripts written in Lua.
			\item Advisor: \href{https://www.antexel.com/sylefeb/research}{Sylvain Lefebvre}
		\end{itemize}}
\end{cvtimeline}

% ── Publications ────────────────────────────────────────────
\section{\faBookOpen\enspace Publications}

\subsection*{Journal articles}

\begin{cvtimeline}
	\cvlastentry{2026}
	{Towards a verified compiler for Distributed PlusCal (extended)}
	{\textbf{Ghilain Bergeron}, Horatiu Cirstea, Stephan Merz}
	{}
	{\textit{Submitted} to Journal of Logical and Algebraic Methods in Programming}
\end{cvtimeline}

\subsection*{Conference and workshop papers}

\begin{cvtimeline}
	\cventry{2025}
	{Formalizing Splitting in Isabelle/HOL \href{https://doi.org/10.4230/LIPIcs.ITP.2025.22}{\textcolor{cvblue}{\faLink}}}
	{\textbf{Ghilain Bergeron}, Florent Krasnopol, Sophie Tourret}
	{}
	{16th International Conference on Interactive Theorem Proving (ITP)}

	\cventry{2025}
	{Towards a verified compiler for Distributed PlusCal \href{https://hal.science/hal-05329156/}{\textcolor{cvblue}{\faLink}}}
	{\textbf{Ghilain Bergeron}, Horatiu Cirstea, Stephan Merz}
	{}
	{11th International Workshop on Rewriting Techniques for Program Transformations and Evaluation (WPTE)}

	\cvlastentry{2024}
	{HOL Light to Isabelle/HOL translation, rebooted \href{https://members.loria.fr/sophie.tourret/papers/isabelle24translation.pdf}{\textcolor{cvblue}{\faLink}}}
	{\textbf{Ghilain Bergeron}, Stéphane Glondu, Sophie Tourret}
	{}
	{Isabelle Workshop}
\end{cvtimeline}

\subsection*{Preprints / submitted}

\begin{cvtimeline}
	\cventry{2026}
	{Encoding Network PlusCal into the Join Calculus}
	{\textbf{Ghilain Bergeron}}
	{}
	{\textit{Submitted}}

	\cvlastentry{2026}
	{BARReL: A modern backend for Atelier B in Lean}
	{\textbf{Ghilain Bergeron}, Vincent Trélat}
	{}
	{\textit{Submitted}}
\end{cvtimeline}

% ── Talks \& presentations ───────────────────────────────────
\section{\faMicrophone\enspace Talks \& Presentations}

\begin{cvtimeline}
	\cventry{Nov 2025}
	{Sémantique Dénotationnelle, espaces métriques et Go}
	{}
	{}
	{Journées LVP du GDR SciLog du CNRS}

	\cventry{Sep 2025}
	{Towards a verified compiler for Distributed PlusCal}
	{}
	{}
	{VeriDis retreat}

	\cventry{Jul 2025}
	{Towards a verified compiler for Distributed PlusCal}
	{}
	{}
	{11th International Workshop on Rewriting Techniques for Program Transformations and Evaluation (WPTE)}

	\cventry{Nov 2024}
	{Compilation vérifiée de spécifications d'algorithmes distribués}
	{}
	{}
	{Journées LVP du GDR SciLog du CNRS}

	\cventry{Aug 2024}
	{An encoding of Distributed PlusCal in the Join Calculus}
	{}
	{}
	{VINO'24}

	\cvlastentry{Jun 2024}
	{An encoding of Distributed PlusCal in the Join Calculus}
	{}
	{}
	{VeriDis retreat}
\end{cvtimeline}


% ── Grants \& Awards ────────────────────────────────────
\section{\faTrophy\enspace Grants \& Awards}

\begin{cvtimeline}
	\cvlastentry{2022 -- 2023}
	{ORION Excellence Scholarship}
	{\href{https://www.univ-lorraine.fr/osezlarecherche/les-dispositifs-orion/bourses-dexcellence/}{University of Lorraine --- ORION Programme}}
	{}
	{}
\end{cvtimeline}

% ── Software / Personal projects ────────────────────────
\section{\faLaptopCode\enspace Software \& Personal Projects}

\begin{cvtimeline}
	\cventry{Oct 2023 -- present}
	{Fugue \textnormal{\small\href{https://github.com/Mesabloo/fugue}{\textcolor{cvblue}{\faGithub}}}}
	{\cvtag{Lean 4}\cvtag{Go}}
	{}
	{A compiler for Distributed PlusCal specifications targeting native Go code,
		with a formally verified compilation core developed in the Lean~4 proof assistant.}

	\cventry{2024}
	{BARReL \textnormal{\small\href{https://github.com/vtrelat/barrel}{\textcolor{cvblue}{\faGithub}}} --- with Vincent Trélat}
	{\cvtag{Lean 4}}
	{}
	{Bridges Atelier B proof obligations to Lean~4 by parsing \texttt{.pog} files,
		encoding B terms and types into Lean expressions, and discharging proof obligations
		interactively using Lean tactics and Mathlib.}

	\cventry{2021 -- 2022}
	{Zilch \textnormal{\small\href{https://github.com/zilch-lang/zilch}{\textcolor{cvblue}{\faGithub}}}}
	{\cvtag{Haskell}\cvtag{Agda}\cvtag{Isabelle/HOL}}
	{}
	{A personal research experiment investigating the combination of dependent types, algebraic effects,
		and region-based memory management in a statically-typed, low-level functional programming language.
		Includes a branch with a typechecker formalised in Agda, and related formal developments in Isabelle/HOL.}

	\cvlastentry{2021 -- 2022}
	{N\textsuperscript{\textbf{*}} \textnormal{\small\href{https://github.com/zilch-lang/nstar}{\textcolor{cvblue}{\faGithub}}}}
	{\cvtag{Haskell}\cvtag{C}}
	{}
	{A personal research project exploring a typed assembly language inspired by FunTAL and related approaches,
		targeting cross-platform compilation with strong type-safety guarantees at the assembly level.}
\end{cvtimeline}

% ── Collective tasks ────────────────────────────────────────
\section{\faUsers\enspace Collective Tasks}

\begin{cvtimeline}
	\cvlastentry{2024}
	{Local organizer}
	{12th International Joint Conference on Automated Reasoning (IJCAR) --- Nancy, France}
	{}
	{}
\end{cvtimeline}

% ── Reviews ────────────────────────────────────────────────
\section{\faSearch\enspace Reviews}

\subsection*{Subreviewer}

\begin{cvtimeline}
	\cventry{2026}
	{Theoretical Aspects of Software Engineering Conference (TASE)}
	{}{}{}

	\cvlastentry{2026}
	{Journées Francophones des Langages Applicatifs (JFLA)}
	{}{}{}
\end{cvtimeline}

% ── Teaching ────────────────────────────────────────────────
\section{\faChalkboardTeacher\enspace Teaching}

\begin{cvtimeline}
	\cventry{2023 -- 2025}
	{Teaching assistant --- 2\textsuperscript{nd} year of master's degree in computer science}
	{University of Lorraine --- Télécom Nancy}
	{}
	{\begin{itemize}
			\item Functional programming course in Haskell.
		\end{itemize}}

	\cvlastentry{2024 -- 2025}
	{Internship supervision --- 1\textsuperscript{st} year of master's degree in computer science}
	{University of Lorraine --- \'Ecole des Mines de Nancy}
	{}
	{\begin{itemize}
			\item Thomas AROCENA \newline
			      Writing and mechanizing an operational semantics for eBPF in the Lean 4 proof assistant.
		\end{itemize}}
\end{cvtimeline}

% % ── References ──────────────────────────────────────────────
% \section{References}

% Available upon request.

\end{document}
